\section{From Fixing Issues to Building Applications}\label{sec:building}

Coding-agent research grew up around issue resolution, where an executed test suite decides whether an agent's patch fixes a reported bug. Headline numbers have risen from 1.96\% of issues for the best model on the original SWE-bench \citep{C1} to 99.4\% for one frontier agent on a repaired split of a harder successor, as reported by that benchmark's owner \citep{C12}. The trend is real, but the instrument has repeatedly proved weaker than assumed. A manual audit of one leading agent's passed patches (SWE-agent with GPT-4) found the solution in the issue report or its comments for 32.67\% of them, and weak tests behind a further 31.08\% \citep{C5}. Differential testing found that 29.6\% of plausible patches behave differently from the developers' fix; together with patches that fail the developers' own tests, these weaknesses inflate reported resolution rates by 6.4 points \citep{V12}. OpenAI stopped reporting SWE-bench Verified after finding evidence of training contamination and flawed tests in at least 59.4\% of a 138-problem subset chosen because its model did not consistently solve those problems \citep{C15}, and the owner of SWE-Bench Pro repaired or removed 480 of its 731 tasks \citep{C12}. Contamination adds to this: the same agent and model scored 43.20\% on SWE-bench Verified but 19.25\% on freshly collected issues \citep{C17}, a gap that mixes memorization with differences in difficulty. And passing the grader is not passing review. In a research note in which 4 maintainers reviewed 296 agent-written pull requests, merge decisions ran about 24 percentage points below the automated grader's scores \citep{C14}.

Building from scratch, from a specification and an empty workspace, is far less solved, and the benchmarks that test it are young. On Vibe Code Bench, whose 100 web-application specifications are graded by a browser agent through 964 workflows, the best of 16 models passes 61.8\% of workflows, and average accuracy falls from 34.18\% for applications without external integrations to 13.49\% for those needing both email and payments \citep{C19}. On NL2Repo-Bench, where one requirements document must become a Python library graded by the real library's tests, every model averages below 40.5\% of tests passed, and the strongest fully passes the suite for only five of 104 repositories in a single run \citep{C22}. Android Bench, on 30 long mobile tasks written by Android developers, puts the best pass rate with hidden tests at 41\% although the mean partial score is 0.85; most failed runs miss a requirement rather than fail to build \citep{C29}. On SWE-Lancer, built from 1,488 paid freelance tasks, the best model in OpenAI's 2025 study (Claude 3.5 Sonnet) earned \$208,050 of the \$500,800 available on the public Diamond set \citep{C16}. The scores are not comparable across benchmarks, but together they suggest that from-scratch work fails at the last requirement, not the first build.

The levers these studies identify fit the thesis: agents build better when they know how they will be checked and can observe what they built. In Android Bench, moving from a closed variant (acceptance criteria only, tests hidden) to an open one (detailed requirements plus failing tests the agent may run) raised the pass rate of the agent that did best with tests available from 22.7\% to 84.0\%, and pass rates rose 2.6 to 6.7 times across the six agents run on both variants \citep{C29}. Because the open variant changes requirement detail as well as test access, the gain cannot be assigned to tests alone. Removing the human-written requirements and interface from SWE-Bench Pro tasks cut GPT-5 from 25.9\% to 8.4\% \citep{C9}: a specification that states what the grader checks is part of the signal. Self-observation tracks success. On Vibe Code Bench, the number of browser calls an agent makes to test its own application correlates with accuracy at $r = 0.72$ \citep{C19}, although a cross-model correlation cannot separate ``testing helps'' from ``stronger models test more''. Harness design for long runs adds a separate checker. In an engineering post by Anthropic about its own models, a planner, a generator and an evaluator that clicks through the running application produced output the author judged clearly better than a solo run, at more than 20 times the cost (\$200 and 6 hours against \$9 and 20 minutes); the author credits ``separating the agent doing the work from the agent judging it'' \citep{T28}. With a single reported comparison, the author as judge and a solo arm that cost a twentieth as much, it suggests an effect but does not establish one.

From-scratch tasks have their own leakage route: code removed from a real repository may still be present. The commit0 benchmark strips function bodies from 57 Python libraries and grades re-implementations against each library's own tests. An issue, since closed as completed, reported that agents could read the original implementations from git history and paste them in (\repo{commit-0/commit0}{113}), and a pull request to exclude test-file changes from the graded diff was closed unmerged after a maintainer noted it was not yet correct (\repo{commit-0/commit0}{100}). The owner of SWE-Bench Pro found the pattern at scale: with the network open, 32 of 642 runs of one model contacted code hosts, and a scripted agent that edits no source file but replaces \texttt{python3} with a script printing PASSED scores 1 when the verifier runs inside the agent's sandbox \citep{C12}. The remedy is to prove each task both ways before any agent runs: a reference solution must pass and an empty attempt must fail. Terminal-Bench requires an oracle solution and a no-op agent that fails \citep{C18}. The harbor harness from the Terminal-Bench team requires benchmark adapters to reach 100\% oracle reward (\repofile{harbor-framework/harbor}{adapters/ADAPTER_CONTRIBUTING.md}) and its task-quality rubric states that the harness copies tests into the container only after the agent runs (\repofile{harbor-framework/harbor}{src/harbor/cli/quality_checker/default-rubric.toml}). For builders, the evidence implies that a from-scratch task is only as trustworthy as the proof that it can be passed, cannot be passed by doing nothing, and cannot be passed by reading the answer.
